\chapter{Performance}
\label{ch:performance}

Yap's workload is small per user but frequent: every live user refreshes a nearby list every 20 seconds,
and every list shows photos. The design keeps each refresh to one database query and as few photo
downloads as possible. \Cref{fig:nearby-path} follows one refresh and notes the work saved at each step;
the rest of the chapter lists the choices that matter, grouped by where they act.

\begin{figure}[ht]
  \centering
  % The nearby refresh path and the work saved at each step (Figure: nearby path).
\begin{tikzpicture}[
    stepbox/.style={box, font=\sffamily\scriptsize, text width=31mm, minimum height=10mm, align=center},
    saving/.style={font=\sffamily\tiny, text=accentdark, text width=33mm, align=center, anchor=north},
  ]
  \node[stepbox, clientbox] (poll) at (1.85,0) {\textbf{One shared poll}\\every 20\,s};
  \node[stepbox, serverbox] (rpc) at (5.95,0) {\textbf{\texttt{nearby()}}\\one query};
  \node[stepbox, accentbox] (paths) at (10.05,0) {\textbf{\texttt{toPeople}}\\picture and small-copy paths};
  \node[stepbox, accentbox] (cache) at (14.15,0) {\textbf{Signed-URL cache}\\\texttt{signPaths}};

  \node[stepbox, serverbox] (sign) at (14.15,-2.9) {\textbf{\texttt{createSignedUrls}}\\one batch for the misses};
  \node[stepbox, clientbox] (state) at (10.05,-2.9) {\textbf{\texttt{state.nearby}}};
  \node[stepbox, clientbox] (avatar) at (5.95,-2.9) {\textbf{Avatar}\\small copy at 64\,pt or less};
  \node[stepbox, clientbox] (img) at (1.85,-2.9) {\textbf{\texttt{expo-image} cache}\\key without the token};

  \draw[flow] (poll) -- (rpc);
  \draw[flow] (rpc) -- (paths);
  \draw[flow] (paths) -- (cache);
  \draw[flow] (cache) -- node[edgelabel, right=2pt, fill=none] {misses} (sign);
  \draw[flow] (sign) -- (state);
  \draw[flow] (cache.south west) ++(0.3,0) -- node[edgelabel, fill=none, sloped, above=1pt] {hits} (state.north east);
  \draw[flow] (state) -- (avatar);
  \draw[flow] (avatar) -- (img);

  \node[saving] at (poll.south) {paused in the background and when not live};
  \node[saving] at (rpc.south) {GiST index; at most 50 cards; no positions};
  \node[saving] at (paths.south) {further photos wait until a profile opens};
  \node[saving] at ($(sign.south)+(0,0)$) {URLs last 1\,h, reused until 10\,min remain};
  \node[saving] at (state.south) {an older refresh never replaces a newer one};
  \node[saving] at (avatar.south) {320\,px file instead of 1080\,px};
  \node[saving] at (img.south) {downloaded once, however often re-signed};
\end{tikzpicture}

  \caption{One nearby refresh. Most refreshes end with cache hits: no new signed URLs and no downloads.}
  \label{fig:nearby-path}
\end{figure}

\section{Database}

\begin{description}
  \item[A table of live users only.] \code{availability} holds a row only while its user is live, and
    rows are deleted on going offline and by the minute job. The table that every nearby query scans is
    therefore as small as the number of people live at that moment, whatever the number of accounts.
  \item[Spatial indexes.] \code{availability.location} and \code{events.location} have GiST indexes, so
    the distance tests in \code{nearby}, \code{near_each_other} and \code{events_near} use an index rather
    than a scan.
  \item[Partial indexes.] Indexes restricted to \code{pending} and \code{accepted} requests stay small,
    because at any moment few requests are waiting or open. They serve the one-request and one-chat checks,
    \code{my_status} and the minute job.
  \item[Indexed rate limits.] Each rate limit counts a user's recent rows. The counts for requests,
    messages and location changes use indexes on the user and the time; the count of events created uses
    the index on the host.
  \item[Bounded results.] \code{nearby} and \code{events_near} return at most 50 rows; chat history is the
    newest 200 messages, read through an index on the chat and time.
  \item[Work done in one call.] Person cards and event cards are assembled in the database, so a list
    needs one round trip. \code{my_status} restores the live session, the chat and the waiting request in
    one call. The send functions return the saved message, so the sender does not read it back.
\end{description}

\section{Network}

\begin{description}
  \item[One shared nearby poll.] A single poll in the \code{(app)} layout feeds the list, the person sheet
    and the prompts. It runs only while the user is live and the app is in the foreground.
  \item[Photos signed in batches and cached.] All the photo paths a screen needs are signed in one
    request. Signed URLs are reused until ten minutes before they expire, so the same photo keeps the same
    address across refreshes.
  \item[Downloads cached by stable keys.] \code{expo-image} caches each photo under its URL without the
    signing token, so a photo is downloaded once even when its URL is signed again.
  \item[Small copies for small avatars.] Each profile picture has a 320-pixel copy. Avatars of 64 points
    or less, such as those in lists, prompts and the chat header, use it, and so does the host's photo
    on an event.
  \item[Lazy further photos.] Lists sign only profile pictures. A person's further photos are signed when
    their profile is opened, keyed by their paths so that a nearby refresh does not sign them again.
  \item[Smaller uploads.] Photos are re-encoded on the phone as JPEG at most 1080 pixels wide, at quality
    0.8, and never enlarged. The resize reuses the image already decoded instead of reading the file again.
  \item[Filtered Realtime.] Each channel filters on one column (\code{to_user}, \code{id},
    \code{request_id} or \code{event_id}), so the server sends a client only the rows it subscribed to.
\end{description}

\section{Rendering}

\begin{description}
  \item[Selective re-rendering.] Components subscribe to slices of the store through selectors, so a
    change re-renders only the components whose slice changed. The React Compiler, enabled as an Expo
    experiment, memoises components automatically.
  \item[One clock for all countdowns.] A single one-second interval drives every countdown on screen and
    stops when none is shown.
  \item[Virtualised lists.] The nearby list, the blocked list and both chat threads use \code{FlatList};
    chat threads are inverted, so they open at the newest message without scrolling.
  \item[Static map markers.] Event pins are custom views drawn once (\code{tracksViewChanges} is off), so
    the map does not redraw them on every frame.
  \item[Stable list data.] The nearby refresh replaces the list only when an answer is current, and event
    reloads keep the messages already loaded for each event.
\end{description}
